
\chapter{سیستم پیشنهاددهی الگوریتمی}

\section{هدف}

Recommendation باید بدون LLM بتواند:
\begin{itemize}
  \item Feed اولیه را بر اساس Interest بسازد؛
  \item با رفتار Save/Spark تطبیق پیدا کند؛
  \item Creator follow را لحاظ کند؛
  \item محتواهای مشابه را پیدا کند؛
  \item Cross-format discovery بدهد؛
  \item دلیل پیشنهاد را شفاف کند.
\end{itemize}

\section{Cold Start}

User جدید Behavior ندارد. پس Signal اول:
\[
\text{Declared Interests}
\]

مثال:
\begin{latin}
\begin{lstlisting}
User:
  Science
  Business
  History
\end{lstlisting}
\end{latin}

Ideaهای Science/Business/History Score اولیه بالاتر می‌گیرند.

\section{Signalهای اصلی}

\begin{itemize}
  \item explicit interest match؛
  \item saved-topic affinity؛
  \item sparked-topic affinity؛
  \item followed creator؛
  \item popularity؛
  \item freshness؛
  \item content-language preference؛
  \item already-seen penalty.
\end{itemize}

\section{فرمول پایه}

\[
Score(i,u) =
w_1 I +
w_2 S +
w_3 K +
w_4 F +
w_5 P +
w_6 R +
w_7 L -
w_8 V
\]

که در آن:
\begin{itemize}
  \item \(I\): Interest match؛
  \item \(S\): Save affinity؛
  \item \(K\): Spark affinity؛
  \item \(F\): Followed creator signal؛
  \item \(P\): Popularity؛
  \item \(R\): Recency/Freshness؛
  \item \(L\): Language fit؛
  \item \(V\): Viewed/consumed penalty.
\end{itemize}

وزن اولیه پیشنهادی:
\begin{longtable}{p{6.8cm}r}
\toprule
Signal & Weight فرضی \\
\midrule
Explicit interest & +5 \\
Saved topic affinity & +4 \\
Sparked topic affinity & +3 \\
Followed creator & +3 \\
Preferred content language & +2 \\
Popularity & 0 تا +1 \\
Freshness & 0 تا +1 \\
Already consumed & penalty \\
\bottomrule
\end{longtable}

\begin{warningbox}
این Weightها Requirement منبع نیستند. Hypothesis فنی هستند و باید با داده آزمایشی Tune شوند.
\end{warningbox}

\section{جلوگیری از رشد بی‌نهایت Affinity}

اگر User صدها Save روی یک Topic داشت، Score نباید خطی بی‌نهایت رشد کند.

مثال ساده:
\[
Affinity = \min(count, cap)
\]

یا Normalization به بازه مشخص.

\section{Jaccard Similarity}

برای Related Ideas:
\[
J(A,B)=\frac{|A\cap B|}{|A\cup B|}
\]

مثال:
\begin{align*}
A &= \{Productivity, Focus, Psychology\}\\
B &= \{Productivity, Focus, Habits\}\\
J(A,B) &= \frac{2}{4}=0.5
\end{align*}

مزیت:
\begin{itemize}
  \item deterministic؛
  \item قابل توضیح؛
  \item بدون Model/Embedding؛
  \item مناسب دفاع دانشگاهی.
\end{itemize}

\section{Cross-format Diversity}

اگر User Idea کتابی درباره Focus می‌بیند، Related list می‌تواند Podcast با Topic مشابه را Boost کند.

مفهوم:
\[
FinalRelatedScore = TopicSimilarity + DiversityBonus
\]

Format متفاوت نباید خودکار بهتر باشد؛ فقط Bonus کوچک برای جلوگیری از monotony.

\section{Popularity}

Popularity نباید صرفاً total sparks باشد چون محتوای قدیمی همیشه برنده می‌شود.

می‌توانیم از Normalized popularity استفاده کنیم:
\begin{itemize}
  \item sparks در بازه اخیر؛
  \item comments؛
  \item save ratio؛
  \item capped score.
\end{itemize}

برای MVP ساده نگه می‌داریم.

\section{Freshness}

Freshness می‌تواند تابعی از age باشد:
\[
Freshness = \max(0, 1 - \frac{AgeDays}{Window})
\]

مثلاً Window = 30 روز برای یک Bonus کوچک.

\section{Recommendation Explanation}

از Signal واقعی Reason می‌سازیم:
\begin{itemize}
  \item «چون Science را انتخاب کرده‌ای.»
  \item «چون چند Idea درباره Focus ذخیره کرده‌ای.»
  \item «چون این Creator را دنبال می‌کنی.»
  \item «موضوع مشابه، اما این بار از یک Podcast.»
\end{itemize}

Reason باید با Score واقعی سازگار باشد.

\section{Feed Ranking و Related Idea جدا هستند}

\subsection{Feed}

User-specific.

\subsection{Related}

Idea-specific و کمتر وابسته به User.

این دو Service جدا طراحی می‌شوند.

\section{ساختار کد پیشنهادی}

\begin{latin}
\begin{lstlisting}
apps/recommendations/
├── scoring.py
├── selectors.py
├── explanations.py
├── services.py
└── tests/
\end{lstlisting}
\end{latin}

\section{Pure Function}

مثال مفهومی:
\begin{latin}
\begin{lstlisting}[language=Python]
def weighted_score(
    interest_match: float,
    saved_affinity: float,
    sparked_affinity: float,
    followed_creator: bool,
    popularity: float,
    freshness: float,
) -> float:
    score = 0.0
    score += 5 * interest_match
    score += 4 * saved_affinity
    score += 3 * sparked_affinity
    score += 3 if followed_creator else 0
    score += popularity
    score += freshness
    return score
\end{lstlisting}
\end{latin}

در Implementation واقعی Inputها دقیق‌تر تعریف می‌شوند.

\section{Performance}

اشتباه:
\begin{latin}
\begin{lstlisting}[language=Python]
ideas = Idea.objects.all()
# calculate every idea in Python for every request
\end{lstlisting}
\end{latin}

بهتر:
\begin{enumerate}
  \item Candidateها را با DB محدود کن؛
  \item language/status/topic filter؛
  \item exclude hidden/blocked؛
  \item candidate window؛
  \item سپس Score.
\end{enumerate}

\section{Evaluation}

Testهای لازم:
\begin{itemize}
  \item cold start؛
  \item interest ranking؛
  \item save/spark adaptation؛
  \item follow boost؛
  \item language filter؛
  \item seen penalty؛
  \item cross-format diversity؛
  \item deterministic result؛
  \item latency زیر Target.
\end{itemize}

\section{TF-IDF به‌عنوان Stretch}

اگر Topic-only similarity ساده بود:
\begin{itemize}
  \item TF-IDF روی Title/Body؛
  \item Cosine Similarity؛
  \item scikit-learn.
\end{itemize}

این قابلیت Optional است و بدون آن Project کامل محسوب می‌شود.

\section{Conflict با Summarize in 30s}

Story منبع Quick Chip مثال \eng{Summarize in 30s} دارد. چون Generative AI حذف شده:
\begin{itemize}
  \item آن Story یا Deferred می‌شود؛
  \item یا به «Ideaهای کوتاه‌تر» تغییر می‌کند؛
  \item یا Short Excerpt دستی از قبل ذخیره می‌شود.
\end{itemize}

این تغییر Scope باید در گزارش نهایی صریح ذکر شود.
